\documentclass[12pt,oneside]{book}

% =========================================================
% METADATOS
% =========================================================
\newcommand{\universidad}{Universidad de Buenos Aires (UBA)}
\newcommand{\materia}{Derecho Civil}
\newcommand{\titulo}{La Lesión como Vicio de los Actos Jurídicos}
\newcommand{\autor}{Isaac Edgar Camacho Ocampo}
\newcommand{\anio}{2026}

% =========================================================
% PAQUETES
% =========================================================
\usepackage[utf8]{inputenc}
\usepackage[T1]{fontenc}
\usepackage[spanish,es-nodecimaldot]{babel}

\usepackage[
    top=2.5cm,
    bottom=2.5cm,
    left=3cm,
    right=2.5cm,
    headheight=15pt
]{geometry}

\usepackage{amsmath}
\usepackage{amssymb}
\usepackage{mathtools}

\usepackage{graphicx}
\usepackage{float}
\usepackage{caption}
\usepackage{subcaption}

\usepackage{booktabs}
\usepackage{array}
\usepackage{longtable}
\usepackage{tabularx}

\usepackage{enumitem}

\usepackage{xcolor}
\usepackage[most]{tcolorbox}

\usepackage{fancyhdr}
\usepackage{ragged2e}

\graphicspath{
    {./img/}
    {./graficos/}
    {./figuras/}
}

% =========================================================
% CONFIGURACIÓN GENERAL
% =========================================================
\setcounter{tocdepth}{2}

\setlength{\parindent}{0pt}
\setlength{\parskip}{0.5em}

\setlist[itemize]{
    topsep=0.3em,
    itemsep=0.2em
}

\setlist[enumerate]{
    topsep=0.3em,
    itemsep=0.2em
}

\clubpenalty=10000
\widowpenalty=10000

% =========================================================
% ENCABEZADOS Y PIES
% =========================================================
\pagestyle{fancy}
\fancyhf{}

\fancyhead[L]{\nouppercase{\leftmark}}
\fancyhead[R]{\materia}

\fancyfoot[C]{\thepage}

\renewcommand{\headrulewidth}{0.4pt}
\renewcommand{\footrulewidth}{0pt}

\fancypagestyle{plain}{
    \fancyhf{}
    \fancyfoot[C]{\thepage}
    \renewcommand{\headrulewidth}{0pt}
}

% =========================================================
% COMANDOS MATEMÁTICOS
% =========================================================
\newcommand{\abs}[1]{\left|#1\right|}
\newcommand{\norm}[1]{\left\lVert#1\right\rVert}
\newcommand{\vect}[1]{\mathbf{#1}}
\newcommand{\deriv}[2]{\frac{d#1}{d#2}}
\newcommand{\pderiv}[2]{\frac{\partial#1}{\partial#2}}

% =========================================================
% CAJAS ACADÉMICAS
% =========================================================
\newtcolorbox{definition}[1]{
    enhanced,
    breakable,
    title={Definición: #1},
    fonttitle=\bfseries,
    colback=blue!3,
    colframe=blue!50!black,
    boxrule=0.8pt,
    arc=2mm
}

\newtcolorbox{important}{
    enhanced,
    breakable,
    title={Importante},
    fonttitle=\bfseries,
    colback=yellow!8,
    colframe=orange!70!black,
    boxrule=0.8pt,
    arc=2mm
}

\newtcolorbox{example}[1]{
    enhanced,
    breakable,
    title={Ejemplo: #1},
    fonttitle=\bfseries,
    colback=green!4,
    colframe=green!40!black,
    boxrule=0.8pt,
    arc=2mm
}

\newtcolorbox{formula}[1]{
    enhanced,
    breakable,
    title={Fórmula / Regla: #1},
    fonttitle=\bfseries,
    colback=purple!4,
    colframe=purple!50!black,
    boxrule=0.8pt,
    arc=2mm
}

\newtcolorbox{exercise}[1]{
    enhanced,
    breakable,
    title={Ejercicio: #1},
    fonttitle=\bfseries,
    colback=gray!5,
    colframe=gray!60!black,
    boxrule=0.8pt,
    arc=2mm
}

\newtcolorbox{note}{
    enhanced,
    breakable,
    title={Nota},
    fonttitle=\bfseries,
    colback=white,
    colframe=gray!50,
    boxrule=0.6pt,
    arc=2mm
}

% =========================================================
% HIPERVÍNCULOS Y METADATOS PDF
% =========================================================
\usepackage[
    colorlinks=true,
    linkcolor=blue!50!black,
    citecolor=green!40!black,
    urlcolor=blue!60!black,
    pdftitle={\titulo},
    pdfauthor={\autor},
    pdfsubject={Apuntes universitarios},
    bookmarks=true
]{hyperref}

% =========================================================
% DOCUMENTO
% =========================================================
\begin{document}

% ---------------------------------------------------------
% PORTADA
% ---------------------------------------------------------
\begin{titlepage}

\thispagestyle{empty}

\begin{center}

\vspace*{1cm}

{\large \textsc{\universidad}}

\vspace{3cm}

{\Huge \textbf{\materia}}

\vspace{1cm}

{\LARGE \titulo}

\vspace{0.8cm}

\textit{Comprender los fundamentos de cada instituto es el camino más seguro hacia un estudio verdaderamente sólido.}

\vspace{1.2cm}

\rule{0.70\textwidth}{0.5pt}

\vspace{0.5cm}

{\large Apuntes de estudio}

\vspace{0.5cm}

\rule{0.70\textwidth}{0.5pt}

\vfill

{\large \autor}

\vspace{0.3cm}

{\large \anio}

\end{center}

\vfill

\begin{flushleft}
\textit{Este es un modesto aporte a los estudiantes de ``Derecho Civil'' no pretende sustituir el material oficial de la materia.}
\end{flushleft}

\end{titlepage}

% ---------------------------------------------------------
% ÍNDICE
% ---------------------------------------------------------
\tableofcontents

% ===========================================================================
\chapter[Antecedentes históricos y derecho comparado]{Antecedentes históricos y derecho comparado de la lesión}
% ===========================================================================

\section{Los vicios de los actos jurídicos según el Código Civil y Comercial}

El estudio de la lesión se inserta dentro del análisis general de los vicios propios de los actos jurídicos. En el ordenamiento argentino vigente, estos vicios son tres: la \textbf{lesión}, la \textbf{simulación} y el \textbf{fraude}.

A diferencia del tratamiento que recibían los vicios de la voluntad —cada uno regulado en un capítulo propio—, en el Código Civil y Comercial (CCC) la lesión, la simulación y el fraude están tratados dentro del capítulo sexto del título cuarto, distribuidos en tres secciones.

\begin{note}
La lesión no estaba incluida originariamente en el Código de Vélez Sarsfield. Se trata de un instituto que resulta central para comprender el funcionamiento concreto de la justicia conmutativa en el derecho privado patrimonial, razón por la cual merece un desarrollo extenso y diferenciado, con antecedentes históricos propios.
\end{note}

\section{Concepto de lesión: una distinción necesaria}

\begin{definition}{Lesión como vicio de los actos jurídicos}
Cuando se habla de lesión como vicio de los actos jurídicos, \textbf{no} se hace referencia a la lesión física —como en el derecho penal, donde el delito de lesiones se clasifica en graves, levísimas y leves—. Se trata de un concepto jurídico distinto, vinculado a la existencia de una \textbf{desproporción entre las prestaciones} de un acto jurídico.
\end{definition}

En todo acto jurídico —especialmente en la compraventa— existe la expectativa de que las prestaciones sean equivalentes. Este principio se vincula con la denominada \textbf{justicia conmutativa}, entendida como la justicia entre particulares que busca una igualdad aritmética: que lo que una parte da sea igual a lo que recibe.

La pregunta central que atraviesa la construcción histórica del instituto es la siguiente: en principio, se supone que las partes son libres y pactan voluntariamente los términos de su acuerdo, de modo que no habría razón para invalidarlo si no existió error, dolo ni violencia. Sin embargo, cuando aparece una desproporción notable entre las prestaciones, surge el interrogante de en qué punto esa desproporción se vuelve tan grave como para dar lugar a la lesión.

A este elemento objetivo —la desproporción— se suma un segundo elemento, de naturaleza subjetiva: la explotación, por una de las partes, de la inferioridad de la otra, con el fin de obtener una ventaja patrimonial evidentemente desproporcionada.

\section{El derecho romano y el Código Justinianeo}

En el derecho romano existen dos rescriptos atribuidos a los emperadores Diocleciano y Maximino que habrían instituido la figura, aunque se ha discutido si se trata de textos insertados con posterioridad. Lo indudable es que el instituto de la lesión aparece recogido en el Código Justinianeo, la recopilación del ordenamiento jurídico romano elaborada bajo el Emperador Justiniano.

\begin{formula}{Código Justinianeo (Corpus Iuris Civilis)}
``Si tú o tu padre hubiera vendido por menor precio una cosa, es humano que, restituyendo todo el precio a los compradores, recobres el fondo vendido mediando la autoridad del juez, aunque si el comprador lo prefiere, reciban lo que les falta al justo precio. Pero se considera que el precio es menor si no se hubiera recibido ni la mitad del verdadero precio.''
\end{formula}

Este instituto recibió el nombre de \textbf{lesión enorme}. Era concedida al vendedor en la compraventa de inmuebles, y exigía que la desproporción alcanzara una medida determinada: más de la mitad del valor real de la cosa.

La lesión enorme daba lugar a dos acciones posibles:

\begin{itemize}
    \item una acción para recuperar el inmueble (el ``fundo''); o bien
    \item una acción para que el comprador tuviera la opción de pagar el complemento del precio faltante, es decir, un reajuste.
\end{itemize}

\begin{example}{Reajuste del precio en el derecho romano}
Si la cosa vale 100 y se pagó 40, el comprador —a quien le interesa conservar la cosa— puede optar por pagar los 60 que faltan, en lugar de restituir el fondo.
\end{example}

\section{La Edad Media y el precio justo}

Los glosadores, que continuaron comentando y ampliando el derecho romano, extendieron la aplicación de la lesión a la mayoría de los contratos durante la Edad Media. Los canonistas, por su parte, pusieron especial énfasis en la inmoralidad del desequilibrio entre las prestaciones.

Este factor moral resulta históricamente relevante, porque en torno a la lesión siempre ha oscilado la tensión entre una mirada objetiva —centrada en la desproporción del objeto— y un reproche dirigido a quien explota una situación ajena.

En este contexto adquiere especial relevancia el aporte de \textbf{Tomás de Aquino} y su idea del \textbf{precio justo}, un concepto de difícil determinación pero de gran influencia doctrinal: cuando el precio resulta manifiestamente injusto, se configura una posible lesión.

\begin{formula}{Tomás de Aquino — El precio justo}
Aun la ley humana —en referencia al derecho romano— considera ilícito que el vendedor venda una cosa por más de lo que vale, o que el comprador la adquiera por menos de su valor, aunque no lo tiene por ilícito si la diferencia no es excesiva. Se reputa ilícito cuando el engaño supera la mitad del valor de la cosa, en clara referencia a la métrica del Código Justinianeo.
\end{formula}

\section{La legislación española como antecedente del derecho argentino}

La evolución de la figura en las fuentes españolas —antecedente directo del derecho argentino— puede sintetizarse del siguiente modo:

\begin{table}[H]
\centering
\caption{Tratamiento de la lesión en la legislación española}
\begin{tabularx}{\textwidth}{>{\bfseries\raggedright\arraybackslash}p{0.28\textwidth} X}
\toprule
\textbf{Fuente española} & \textbf{Tratamiento de la lesión} \\
\midrule
Fuero Juzgo & Rigió durante la dominación visigoda. No admite la lesión, pero combate seriamente la usura, instituto muy cercano. \\
Fuero Real & Concede la acción al vendedor cuando lo engañan en más de la mitad del valor, si el comprador no acepta pagar la diferencia. \\
Las Leyes de Partidas & Consolidan la acción tanto para el vendedor como para el comprador. \\
Novísima Recopilación & Acepta la acción en compraventa, permuta, rentas y otros contratos semejantes. \\
\bottomrule
\end{tabularx}
\end{table}

\section{El derecho comparado moderno}

\subsection{El Código Civil Francés (1804): la lesión objetiva}

El Código Napoleónico constituye el ejemplo por excelencia de lo que se ha denominado \textbf{lesión objetiva}: aquella que contempla únicamente los elementos objetivos —es decir, la desproporción— sin tomar en consideración, en la configuración de la norma, el elemento subjetivo de la explotación de la inferioridad de la otra parte.

\begin{formula}{Código Civil Francés — art. 1.674 (1804)}
Cuando el vendedor resulte lesionado en más de siete doceavos del precio del inmueble, tendrá derecho a solicitar la rescisión de la venta, aunque en el contrato haya renunciado expresamente a solicitar dicha rescisión y aunque haya declarado donar las plusvalías.
\end{formula}

La medida establecida —más de 7/12— representa un poco más que la mitad (que serían 6/12). El código francés sigue así la métrica del derecho romano, pero la precisa mediante un porcentaje específico, tomando el valor al momento de la venta.

\begin{example}{Complemento del precio en el código francés}
Se concede al comprador —lesionante— la posibilidad de ofrecer el complemento del precio con la deducción de la décima parte. Si la cosa vale 100 y se pagaron 40, el comprador puede ofrecer el complemento hasta 90, pagando los 50 restantes, y de esa manera conservar la propiedad.
\end{example}

En el derecho argentino, esta acción de complemento recibe el nombre de \textbf{reajuste equitativo} y es fijada por el juez. En el código francés, la prescripción de la acción era de dos años.

\subsection{El Código Civil Alemán — BGB (1900): la lesión subjetiva}

En 1900, la reforma del código civil alemán incorpora por primera vez la fórmula de la denominada \textbf{lesión subjetiva}, estrechamente vinculada al combate de la usura. La doctrina coincide en señalar este momento como el surgimiento de dicha categoría.

\begin{formula}{Código Civil Alemán — BGB (síntesis, 1900)}
Un acto jurídico contrario a las buenas costumbres por el cual alguien, explotando la necesidad, la ligereza o la inexperiencia de otro, obtiene para sí o para un tercero que a cambio de una prestación le prometa o entregue ventajas patrimoniales que exceden de tal forma el valor de la prestación que, teniendo en cuenta las circunstancias, existe una desproporción chocante con ella.
\end{formula}

Aparecen aquí los elementos subjetivos del lesionante —que explota la necesidad, la ligereza o la inexperiencia— y se hace prometer o entregar ventajas patrimoniales. Existe, entonces, una acción activa: el lesionante toma la iniciativa para explotar una inferioridad. Del lado del lesionado, la inferioridad se define como necesidad, ligereza o inexperiencia, términos que luego reproducirá casi textualmente el art.\ 954 argentino y, con alguna variante, el código vigente.

\begin{important}
El código alemán declaraba el acto nulo en su totalidad, sin admitir nulidades parciales ni reajuste de precios, por tratarse de un acto contrario a las buenas costumbres. Existía un deber de restituir; la acción no era confirmable, era imprescriptible y podía ser declarada de oficio.
\end{important}

\subsection{El Código Suizo (1912) y el Código Italiano (1942)}

En 1912, el Código Suizo de las Obligaciones sigue al modelo alemán con variantes propias: coloca la prescripción en un año, torna la acción anulable y permite la confirmación del acto, acercándose así a la idea de un vicio que da lugar a una nulidad relativa susceptible de confirmación.

En 1942 se sanciona el Código Civil Italiano, que ejercerá una influencia decisiva sobre la reforma del art.\ 954 argentino en 1968.

\begin{formula}{Código Civil Italiano — art. 1.448 (1942, síntesis)}
Si hubiera desproporción en las prestaciones de una de las partes, y esa desproporción dependiese del estado de una de ellas, del cual se aprovechara la otra para obtener ventajas, la parte damnificada podrá demandar la rescisión del contrato. La acción no es admisible si la lesión es inferior a la mitad del valor que la prestación ejecutada o prometida por la parte perjudicada tenía en el momento del contrato. La lesión debe perdurar hasta el momento en que se proponga la demanda.
\end{formula}

\begin{note}
La exigencia de que la lesión perdure hasta el momento de la demanda no estaba presente en los códigos anteriores, y será recogida tanto por la ley 17.711 como por el CCC vigente. En el código italiano, además, los contratos aleatorios no pueden ser rescindidos por esta vía.
\end{note}

% ---------------------------------------------------------
% CORNELL — CAPÍTULO 1
% ---------------------------------------------------------
\section*{Cuadro de repaso — Método Cornell}
\addcontentsline{toc}{section}{Cuadro de repaso — Método Cornell}

\begin{table}[H]
\centering
\begin{tabularx}{\textwidth}{
    >{\raggedright\arraybackslash}p{0.30\textwidth}
    >{\raggedright\arraybackslash}X
}
\toprule
\textbf{Preguntas / palabras clave} & \textbf{Notas / respuestas} \\
\midrule

\textbf{¿Qué son los vicios propios de los actos jurídicos en el CCC?} &
Son tres: lesión, simulación y fraude. Están regulados en el capítulo sexto del título cuarto del Código Civil y Comercial, distribuidos en tres secciones. \\

\textbf{¿Qué es la lesión como vicio jurídico?} &
No es la lesión física, sino una desproporción entre las prestaciones de un acto jurídico, combinada con la explotación, por una de las partes, de la inferioridad de la otra. Se vincula con la idea de justicia conmutativa (igualdad aritmética entre lo que se da y lo que se recibe). \\

\textbf{¿Cuál es el origen histórico del instituto?} &
El Código Justinianeo (siglo VI) consagró la ``lesión enorme'' en la compraventa de inmuebles: si el precio era inferior a la mitad del valor real, el vendedor podía pedir la nulidad o el complemento del precio. \\

\textbf{¿Qué distinción clave introduce el BGB alemán de 1900?} &
Introduce la lesión subjetiva: ya no basta la desproporción objetiva, sino que se requiere además la explotación de la necesidad, ligereza o inexperiencia de la otra parte. Esta fórmula fue luego adoptada por Argentina en 1968 y rige, con variantes, en el art.\ 332 del CCC. \\

\midrule

\multicolumn{2}{p{\dimexpr\textwidth-2\tabcolsep\relax}}{
\textbf{Resumen:} La lesión, desde la lesión enorme del Código Justinianeo (más de la mitad del valor en la compraventa de inmuebles), pasando por el énfasis moral de la Edad Media y el precio justo de Tomás de Aquino, evolucionó hacia dos grandes modelos de derecho comparado: la \textbf{lesión objetiva} del código francés (que exige únicamente la desproporción) y la \textbf{lesión subjetiva} del BGB alemán (que agrega la explotación de la inferioridad de la otra parte), modelo este último que —junto con el código italiano de 1942— influirá decisivamente en la incorporación de la figura al derecho argentino.
} \\

\bottomrule
\end{tabularx}
\end{table}

% ===========================================================================
\chapter{La lesión en el derecho argentino}
% ===========================================================================

\section{El Código de Vélez (1869): la exclusión de la lesión}

Vélez Sarsfield no incluyó la lesión entre los vicios de los actos jurídicos. Al tratar los actos jurídicos en el art.\ 944, y luego de referirse a la simulación, el codificador dedicó una extensa nota —ubicada luego del art.\ 943— para explicar por qué no adoptaba la figura.

\begin{note}
Técnicamente, la nota no se refiere al contenido del art.\ 943 —que trataba sobre la violencia—, sino que Vélez la ubicó allí, antes del comienzo del tratamiento de los actos jurídicos, con el propósito de explicar en algún lugar del código por qué no incorporaba la lesión. Años más tarde, Borda señaló en un fallo de 1958 que una nota al pie no constituye ley.
\end{note}

Vélez ofrece dos grandes razones para no reconocer la lesión:

\begin{formula}{Nota de Vélez al art. 943 — Primer argumento}
``Para sostener nosotros que la lesión enorme, menor o mínima vicia los actos y obtenernos por tanto de proyectar disposiciones de la materia, basta comparar las diversas legislaciones y de las diferencias entre ellas resultados que no han tenido un principio uniforme al establecer esta ley.''
\end{formula}

El codificador observa que la legislación comparada presenta numerosas diferencias: unas la conceden al vendedor, otras al comprador; unas la aplican a la compraventa, otras a distintos actos; unas imponen una medida de más de la mitad, otras no. Ante esta disparidad de criterios, Vélez concluye que no encuentra un principio uniforme y que, por lo tanto, no se siente obligado a incorporar la figura.

\begin{formula}{Nota de Vélez al art. 943 — Segundo argumento}
``Dejaríamos de ser responsables de nuestras acciones si la ley nos permitiera enmendar todos nuestros errores o todas las consecuencias de ellos. El consentimiento libre prestado sin error, dolo y violencia y con la solemnidad requerida por las leyes debe hacer irrevocables los contratos.''
\end{formula}

Este segundo argumento expresa que lo pactado sin error, dolo ni violencia es válido, y que quien se equivoca debe hacerse cargo de las consecuencias de su propio error. Se advierte aquí una visión marcadamente individualista y liberal, coincidente con otras opciones adoptadas por Vélez a lo largo de todo el código.

\section{La jurisprudencia anterior a la ley 17.711}

Durante el siglo XX, distintos fallos judiciales fueron invalidando actos manifiestamente desproporcionados, invocando generalmente el art.\ 953 —referido al objeto del acto—, que disponía que el objeto no podía consistir en cosas o hechos contrarios a la moral. Cuando se pactaba algo contrario a la moral, ello podía dar lugar a la nulidad, de modo que el vicio de lesión se introducía por vía del art.\ 953.

La nota de Vélez al art.\ 943 seguía vigente, pero algunos jueces —entre ellos, Borda, en un fallo de la Cámara Nacional de Apelaciones de 1958 conocido como el caso ``Orlando''— sostuvieron que las notas no constituyen texto legal y que la lesión resulta incompatible con las buenas costumbres.

\begin{example}{Casos de aplicación jurisprudencial de la lesión (antes de 1968)}
Préstamos usurarios, honorarios exorbitantes del administrador de una sucesión, venta a precio vil (muy bajo) de un terreno, y moderación de cláusulas penales excesivas.
\end{example}

En 1961, el Tercer Congreso de Derecho Civil tuvo una importancia decisiva para la reforma de la ley 17.711, al formular recomendaciones que constituyeron, en gran medida, la base de la reforma de 1968.

\section{La ley 17.711 y la reforma al art. 954 (1968)}

La ley 17.711 incorporó diversas figuras al Código Civil —entre ellas, el abuso del derecho y modificaciones en los efectos de la declaración de voluntad en el tiempo— y también la lesión, incorporada en el art.\ 954.

\begin{formula}{Art. 954 — Ley 17.711 (1968, texto incorporado)}
``También podrá demandarse la nulidad o modificación de los actos jurídicos cuando una de las partes, explotando la necesidad, ligereza o inexperiencia de la otra, obtuviera por medio de ellos una ventaja patrimonial evidentemente desproporcionada y sin justificación. Se presume, salvo prueba en contrario, que existe tal explotación en caso de notable desproporción de las prestaciones. Los cálculos deben hacerse según valores al tiempo del acto y la desproporción deberá subsistir al momento de la demanda. Solo el lesionado o sus herederos podrán ejercer la acción cuya prescripción operará a los cinco años. El accionante tiene opción para demandar la nulidad o el reajuste equitativo del convenio, pero la primera de estas acciones se transformará en acción de reajuste si éste fuera ofrecido por el demandado.''
\end{formula}

En este texto se incorporan claramente un elemento objetivo —la desproporción— y un elemento subjetivo —la explotación de la necesidad, ligereza o inexperiencia—. El nuevo código sustituirá luego la palabra ``ligereza'' por ``debilidad psíquica''. Cabe destacar que el texto del art.\ 954 no establece una métrica fija para medir el elemento objetivo.

\section{El art. 332 del CCC: texto y diferencias con el art. 954}

El Código Civil y Comercial vigente sigue, en lo sustancial, al art.\ 954 para regular la lesión, en un único artículo ubicado en la sección primera del capítulo sexto del título cuarto.

\begin{important}
Aunque se trate de un único artículo, el art.\ 332 contiene un desarrollo normativo extenso, que exige un análisis detenido de cada uno de sus elementos.
\end{important}

\begin{formula}{Art. 332 — Código Civil y Comercial de la Nación (vigente)}
``Puede demandarse la nulidad o modificación de los actos jurídicos cuando una de las partes, explotando la necesidad, debilidad psíquica o inexperiencia de la otra, obtuviera por medio de ellos una ventaja patrimonial evidentemente desproporcionada y sin justificación. Se presume, salvo prueba en contrario, que existe tal explotación en casos de notable desproporción de las prestaciones. Los cálculos deben hacerse según valores al tiempo del acto y la desproporción debe subsistir al momento de la demanda. El afectado tiene opción para demandar la nulidad o un reajuste equitativo del convenio, pero la primera de estas acciones se transforma en acción de reajuste si éste es ofrecido por el demandado al contestar la demanda. Solo el lesionado o sus herederos podrán ejercer la acción.''
\end{formula}

Las diferencias del CCC respecto del art.\ 954 pueden sintetizarse así:

\begin{table}[H]
\centering
\caption{Diferencias entre el art. 954 (ley 17.711) y el art. 332 (CCC)}
\begin{tabularx}{\textwidth}{>{\bfseries\raggedright\arraybackslash}p{0.20\textwidth} X}
\toprule
\textbf{Modificación} & \textbf{Explicación} \\
\midrule
Cambio 1 & La palabra ``ligereza'' es reemplazada por ``debilidad psíquica''. El nuevo código clarifica así un punto que había dado lugar a un problema interpretativo, entre quienes ubicaban la ligereza como un obrar irreflexivo sin mayores requisitos y quienes la entendían como expresión de una cierta debilidad psíquica; el CCC se pronuncia por esta última postura. \\
Cambio 2 & El plazo de prescripción se reduce de cinco a dos años. Metodológicamente resulta correcto tratar el plazo en el libro dedicado a las prescripciones, y no dentro del propio artículo. \\
\bottomrule
\end{tabularx}
\end{table}

\section{Ámbito de aplicación del art. 332}

El art.\ 332 recoge, en principio, todos los actos jurídicos en general. Sin embargo, el acto debe ser bilateral y no a título gratuito, puesto que en los actos a título gratuito siempre existe una desproporción entre las prestaciones que resulta ajena a la lógica de la lesión.

\section{El elemento objetivo: la desproporción de las prestaciones}

El primer elemento exigido por la norma es el elemento objetivo: la desproporción de las prestaciones. A diferencia del código francés, que establece un criterio fijo (7/12 partes), el CCC habla de ``ventaja evidentemente desproporcionada'' y de ``notable desproporción'', sin fijar un criterio numérico. No obstante, se exige que la diferencia sea muy notable: no bastaría, por ejemplo, una diferencia del veinte por ciento.

Un requisito adicional es que la desproporción debe ser \textbf{sin justificación}.

\begin{example}{Desproporción con justificación}
Un tío que va a asistir al casamiento de un sobrino le vende un departamento que vale 100.000 dólares en 40.000, como regalo de casamiento. Tras su muerte, otros herederos —primos del beneficiado— demandan por lesión, invocando la diferencia entre el valor real y el precio pagado. El demandado se defenderá alegando que existió una justificación: la intención de realizar una liberalidad. Podrá discutirse si ese regalo excede o no la porción disponible —cuestión de legítima—, pero no podrá hablarse de lesión en sentido propio, porque no hubo explotación sino una justificación para la desproporción.
\end{example}

Los cálculos deben realizarse tomando en cuenta los valores al momento del acto, puesto que todo vicio afecta al acto desde su origen. Sin embargo, la desproporción debe además subsistir al momento de la demanda, exigencia que proviene del código italiano.

\begin{example}{Desproporción que no subsiste}
Una persona compra muy barato un inmueble en un barrio cotizado. La desproporción existe al momento del acto, pero luego el valor de la propiedad desciende notablemente porque el barrio se deteriora y se vuelve inhabitable. Como la desproporción no subsiste al momento de la demanda, no existe acción por lesión, aun cuando el vicio se hubiera configurado al momento del acto.
\end{example}

\section{Los elementos subjetivos}

El art.\ 332 exige, además del elemento objetivo, la \textbf{explotación} por parte del lesionante. Se trata de algo más intenso que un simple ``aprovechar'' —entendido como que la otra parte no advierte la desproporción y el lesionante simplemente continúa adelante—: existe aquí una acción más activa, consistente en ``hacerse dar o prometer'' aprovechando la inferioridad de la otra parte. Este elemento resulta, en la práctica, de difícil prueba.

Del lado del lesionado, el CCC exige que se configure alguno de los siguientes tres estados de inferioridad:

\begin{table}[H]
\centering
\caption{Estados de inferioridad del lesionado (art. 332 CCC)}
\begin{tabularx}{\textwidth}{>{\bfseries\raggedright\arraybackslash}p{0.22\textwidth} X}
\toprule
\textbf{Estado de inferioridad} & \textbf{Definición y aclaración} \\
\midrule
Necesidad & Falta de lo indispensable para la vida. No basta con alegar, por ejemplo, que se tenía necesidad de asistir a la final de un mundial y por ello se vendió una propiedad muy barata. La necesidad exige una situación urgente, vital e indispensable —como cubrir una operación médica de un familiar—, no un deseo o conveniencia. \\
Debilidad psíquica & Presencia de alguna alteración mental. Reemplaza a la palabra ``ligereza'' del código anterior, que había generado dos posturas doctrinarias: obrar irreflexivo sin mayores requisitos, o expresión de una cierta debilidad psíquica. El nuevo código se pronuncia por esta última. \\
Inexperiencia & Falta del conocimiento que otorgan la vida y las relaciones. Quien se dedica profesionalmente al mercado inmobiliario durante treinta años no puede alegar inexperiencia en esa materia; distinta es la situación de una persona que se traslada de otro lugar, desconoce el ambiente y carece de formación en la materia del acto. \\
\bottomrule
\end{tabularx}
\end{table}

\begin{important}
En general se interpreta que estos tres supuestos de inferioridad constituyen una enumeración taxativa: no puede invocarse otro distinto. El proyecto de reforma del CCC del año 2018 proponía ampliar estos supuestos, pero en el derecho vigente son únicamente tres: necesidad, debilidad psíquica e inexperiencia.
\end{important}

\section{La presunción legal del art. 332}

El art.\ 332, al igual que su antecedente el art.\ 954, establece que se presume —salvo prueba en contrario— la existencia de explotación en caso de notable desproporción de las prestaciones.

\begin{important}
Esta presunción constituye una inversión de la carga de la prueba: si el lesionado prueba la notable desproporción de las prestaciones, se presume que hubo explotación, y corresponderá entonces al lesionante —demandado— probar que no existió tal explotación. Existe así una preponderancia práctica del elemento objetivo dentro de una norma que, en su configuración, es de lesión subjetiva.
\end{important}

Parte de la doctrina discute si el término ``notable'' equivale a ``evidentemente desproporcionada''. En general, se considera que lo notable es aquello que salta a la vista sin necesidad de una investigación mayor.

Existe además una discusión acerca del alcance de la presunción: una postura sostiene que solo alcanza al elemento subjetivo del lesionante, debiendo probarse por separado la necesidad, debilidad psíquica o inexperiencia del lesionado; otra postura entiende que, probada la desproporción, ambos elementos subjetivos se presumen.

\begin{note}
Esta presunción legal constituye una novedad propia del derecho argentino, sin equivalente en el código italiano, el código suizo ni el código alemán, y es la que otorga al elemento objetivo su especial preponderancia práctica.
\end{note}

\section{Las acciones, la legitimación y la prescripción}

La lesión da lugar a dos acciones:

\begin{enumerate}
    \item \textbf{Acción de nulidad}, de carácter relativo —susceptible de confirmación—, en tanto está establecida para la protección de la parte lesionada, al igual que ocurre con los restantes vicios.
    \item \textbf{Reajuste equitativo}, mediante el cual el juez corrige la desproporción en el caso concreto, determinando lo que resulte justo. A diferencia del esquema del código francés, que deducía una décima parte para el pago del complemento del precio, en el derecho argentino el reajuste queda librado a la valoración judicial, sin que la ley fije una pauta.
\end{enumerate}

\begin{important}
El reajuste equitativo prima sobre la acción de nulidad cuando es ofrecido por el demandado al contestar la demanda, en virtud del principio de conservación de los actos jurídicos.
\end{important}

La legitimación para iniciar la acción corresponde exclusivamente al lesionado o a sus herederos, tal como establecía ya el art.\ 954. El plazo de prescripción es de \textbf{dos años}, contado desde la fecha en que la obligación debía cumplirse —plazo que la ley 17.711 fijaba en cinco años.

\section{Figuras afines: cláusula penal y usura}

El art.\ 794, relativo a la cláusula penal, guarda semejanza con la lesión: los jueces pueden reducir las penas establecidas en una cláusula penal —cláusula contractual que fija sanciones para el caso de incumplimiento— cuando resulten excesivas. Se trata de una figura próxima a la lesión, en cuanto también existe desproporción susceptible de reducción judicial, aunque con una diferencia relevante: no produce nulidad, sino únicamente reajuste.

La \textbf{usura}, en cambio, constituye un delito en el derecho penal, regulado por el art.\ 175 bis del Código Penal.

\begin{formula}{Art. 175 bis — Código Penal (síntesis) — Usura}
Quien, aprovechando la necesidad, ligereza o inexperiencia de otro, se hiciera dar o prometer, en cualquier forma, para sí o para un tercero, intereses u otras ventajas pecuniarias evidentemente desproporcionadas con su prestación, u otorgare a sabiendas un crédito usurario, será reprimido.
\end{formula}

En la figura penal se requiere dolo. Debe notarse que esta norma no utiliza el término ``explotación'' sino ``aprovechar'', y exige que el autor se haga dar o prometer las ventajas: existe una acción activa dirigida a lograr esa entrega o promesa, que constituye el núcleo del delito.

\section{Debate doctrinario: la esencia de la lesión}

Existe un debate doctrinario acerca de cuál es la esencia de la lesión. El profesor Tobías, en su Tratado de Derecho Civil, identifica tres posturas posibles:

\begin{table}[H]
\centering
\caption{Posturas doctrinarias sobre la naturaleza jurídica de la lesión}
\begin{tabularx}{\textwidth}{>{\bfseries\raggedright\arraybackslash}p{0.28\textwidth} X}
\toprule
\textbf{Postura} & \textbf{Contenido} \\
\midrule
Vicio del objeto & Enfatiza el elemento objetivo: existe un problema en el objeto del acto, que resulta desproporcionado. Para Tobías, esta postura ignora la presencia de la explotación de la inferioridad. \\
Vicio de la voluntad & Enfatiza el elemento subjetivo del lesionado: por su inferioridad —necesidad, debilidad psíquica o inexperiencia— estaría viciada su voluntad. Tobías objeta que, de seguirse este criterio, resultaría irrelevante la existencia de desproporción. \\
Comportamiento reprochable (Tobías) & Es una respuesta del ordenamiento frente a un comportamiento reprochable, vinculado a la buena fe. Se reprocha la explotación. Esta postura se encuentra más cerca del antecedente alemán. \\
\bottomrule
\end{tabularx}
\end{table}

En definitiva, los tres elementos —objeto, voluntad y conducta— están en juego a la hora de determinar la esencia de la lesión, y de fondo siempre se encuentra la justicia conmutativa. El elemento objetivo resulta indudable: el objeto del acto debe ser proporcionado, y si existe desproporción, el acto no realiza plenamente la justicia. Al mismo tiempo, existe una valoración desde las buenas costumbres, que fundamenta el reproche jurídico cuya consecuencia es la nulidad.

\begin{note}
En síntesis del cierre de la clase: el vicio de lesión en el ordenamiento argentino se recorrió desde el Código de Vélez —que no lo contempló, por las razones expuestas en su nota al art.\ 943— pasando por la jurisprudencia posterior y la reforma introducida por la ley 17.711 en el art.\ 954, cuyo texto es sustancialmente reproducido por el art.\ 332 del CCC vigente. Dentro de esta figura se analizaron su fuente, su método, su ámbito de aplicación, sus elementos, las acciones a las que da lugar, la legitimación activa, el plazo de prescripción, su comparación con otras figuras afines y el problema de su esencia jurídica.
\end{note}

% ---------------------------------------------------------
% CORNELL — CAPÍTULO 2
% ---------------------------------------------------------
\section*{Cuadro de repaso — Método Cornell}
\addcontentsline{toc}{section}{Cuadro de repaso — Método Cornell}

\begin{longtable}{
    >{\raggedright\arraybackslash}p{0.30\textwidth}
    >{\raggedright\arraybackslash}p{0.58\textwidth}
}
\toprule
\textbf{Preguntas / palabras clave} & \textbf{Notas / respuestas} \\
\midrule
\endhead

\textbf{¿Por qué Vélez no incluyó la lesión?} &
Por dos razones: (1) la disparidad de criterios en el derecho comparado, que le impedía identificar un principio uniforme; (2) su visión liberal-individualista, según la cual, si no hubo error, dolo ni violencia, el contrato es válido y cada parte debe hacerse cargo de sus propias decisiones. \\

\textbf{¿Cómo incorporó Argentina la lesión antes de 1968?} &
La jurisprudencia —especialmente el fallo de Borda en el caso ``Orlando'' (1958)— la aplicó por vía del art.\ 953 (objeto contrario a la moral). El Tercer Congreso de Derecho Civil (1961) recomendó luego su incorporación expresa. \\

\textbf{¿Qué exige el art. 332 del CCC para configurar la lesión?} &
Un elemento objetivo: desproporción evidente y sin justificación, que exista al momento del acto y subsista al momento de la demanda. Dos elementos subjetivos: explotación por parte del lesionante, y necesidad, debilidad psíquica o inexperiencia del lesionado. \\

\textbf{¿Cómo opera la presunción del art. 332?} &
Si se prueba la notable desproporción de las prestaciones, se presume —salvo prueba en contrario— la explotación. Esto invierte la carga probatoria: es el demandado quien debe demostrar que no explotó. \\

\textbf{¿Qué acciones otorga la lesión?} &
Nulidad relativa (confirmable) o reajuste equitativo (fijado por el juez). El reajuste prima sobre la nulidad cuando es ofrecido por el demandado al contestar la demanda, en virtud del principio de conservación del acto. \\

\textbf{¿Cuál es el plazo de prescripción?} &
Dos años, contados desde que la obligación debía ser cumplida. El CCC redujo el plazo de cinco años que establecía la ley 17.711. \\

\textbf{¿Cuál es la naturaleza jurídica de la lesión según la doctrina?} &
Existe debate entre tres posturas: (1) vicio del objeto; (2) vicio de la voluntad; (3) conducta reprochable, contraria a la buena fe y a las buenas costumbres. Tobías se inclina por esta tercera postura, más próxima al modelo alemán. \\

\midrule

\multicolumn{2}{p{0.88\textwidth}}{
\textbf{Resumen:} Argentina, luego de que Vélez la excluyera por razones liberales y por la falta de un criterio uniforme en el derecho comparado, incorporó la lesión en 1968 mediante la ley 17.711, siguiendo el modelo alemán e italiano. El CCC vigente la regula en el art.\ 332 con tres elementos: uno objetivo (desproporción evidente, sin justificación, que subsista al tiempo de la demanda), uno subjetivo del lesionante (explotación activa) y uno subjetivo del lesionado (necesidad, debilidad psíquica o inexperiencia). La presunción legal de explotación ante la notable desproporción —novedad argentina sin equivalente en los códigos extranjeros— otorga preponderancia práctica al elemento objetivo. Las acciones disponibles son la nulidad relativa y el reajuste equitativo, prevaleciendo este último por el principio de conservación del acto. El plazo de prescripción es de dos años.
} \\

\bottomrule
\end{longtable}

\end{document}
